% Generated draft: requires mathematical and source review.
\documentclass[11pt]{article}
\usepackage[margin=25mm]{geometry}
\usepackage{fontspec}
\setmainfont{Latin Modern Roman}
\usepackage{amsmath,amssymb,mathtools,unicode-math}
\setmathfont{Latin Modern Math}
\usepackage{enumitem,needspace,xurl,hyperref}
\setlength{\parindent}{0pt}
\setlength{\parskip}{5pt}
\setlength{\emergencystretch}{3em}
\newcommand{\N}{\mathbb{N}}
\newcommand{\Z}{\mathbb{Z}}
\newcommand{\Q}{\mathbb{Q}}
\newcommand{\R}{\mathbb{R}}
\newcommand{\C}{\mathbb{C}}
\newcommand{\F}{\mathbb{F}}
\newcommand{\A}{\mathbb{A}}
\newcommand{\PP}{\mathbb{P}}
\newcommand{\E}{\mathbb{E}}
\newcommand{\eps}{\varepsilon}
\newcommand{\dd}{\,\mathrm{d}}
\begin{document}
\section*{UnsolvedMath ID: 10400098}
\textbf{Source problem code:} catalogue entry
\subsection{Definitions and mathematical statement}
A quandle is a set \(X\) with a binary operation \(*\) such that \(x*x=x\), each right translation \(R_y:x\mapsto x*y\) is bijective, and
\[
 (x*y)*z=(x*z)*(y*z)\qquad(x,y,z\in X).
\]
It is connected when the group generated by the \(R_y\) acts transitively on \(X\). Fix such a quandle \(X\) of finite order \(|X|\). Let \(Q(K)\) be the fundamental knot quandle of a knot \(K\), with the source's orientation and knot-category conventions. In a diagram presentation it is generated by arcs with the usual quandle crossing relations.

A quandle representation in the literal homomorphism sense is a map \(f:Q(K)\to X\) satisfying
\[
 f(a*b)=f(a)*f(b)\qquad(a,b\in Q(K)).
\]
Write \(\operatorname{Hom}_{\rm Qnd}(Q(K),X)\) for these maps and \(N_X(K)=|\operatorname{Hom}_{\rm Qnd}(Q(K),X)|\) for their number. For the usual finitely presented knot quandle this set is finite because the target is finite. Whether the source instead counts equivalence classes of representations, or excludes constant representations, must be checked; neither restriction is stated.

Problem 5.2 asks to describe these representations, or their number, for a fixed finite connected target quandle. It suggests doing so using previously known knot invariants, or reducing the computation to target quandles of smaller order. The latter means target size strictly below \(|X|\), but no class of permitted reductions, compatibility data, or reconstruction formula is supplied. The question does not identify a list of known invariants, a complexity requirement, or a precise desired form of description. These are a program of computation and comparison, not an assertion that a particular unnamed invariant already determines \(N_X(K)\).
\subsection{Short English statement}
For a fixed finite connected quandle, describe the homomorphisms from a knot's fundamental quandle and count them. Seek a relation with known knot invariants or a reduction to smaller target quandles. The source leaves representation equivalence and the required form of such a reduction implicit.
\subsection{Sources}
\begin{itemize}
\item Local UnsolvedMath registry, permanent ID \(10400098\), exact source code \texttt{catalogue entry}, Problem 5.2. No individual proposer or proposed year is supplied; the title is truncated.
\item Ohtsuki, \emph{Problems on invariants of knots and 3-manifolds} (2002), Problem 5.2, PDF page 86: \url{https://msp.org/gtm/2002/04/gtm-2002-04-024s.pdf}. Bold headings distinguish 231 numbered items from references in surrounding exposition.
\item Extraction pdf-font-structure, accessed 2026-07-29. Source evidence NEEDS\_REVIEW presents an open problem and requires a current primary-literature audit. Supplied classification OPEN-TRIAGE calls the fixed-connected-quandle representation question not obviously resolved, without further specific references.
\item Dataset attribution: Ulam AI / UnsolvedMath Contributors, CC BY 4.0. Source rights NEEDS\_REVIEW: publicly accessible author/publisher copy, redistribution terms requiring review.
\end{itemize}
\end{document}
